\subsection{TAKE OVER}\label{take-over}

\subsubsection{interactive in progress / extinction in
progress}\label{interactive-in-progress-extinction-in-progress}

There is a church, except there isn't, because the thing we call a
church turns out to contain at least two incompatible objects: the stone
container and the people who once made the container mean something, and
in these photographs the first has survived the second, which is already
suspicious, because normally we speak as though institutions are the
durable things and people are the passing things, whereas here the
institution has literally become rubble and the people are absent and
yet somehow the absence has more pressure than the masonry.

So: church minus ecclesia.

A shell after its assembly.

A room after the reason for the room has left.

Good starting condition.

TAKE OVER begins here, among abandoned churches, holes, tombstones, bits
of marble, damaged geometry, light behaving like it knows more than we
do, and a person who appears with a black square whose status is,
initially, unclear. Object? Token? Signal? Error? Device? Offering? It
is shown to the outside world. It is carried inward. It is placed
somewhere. It reappears. At some point it acquires the visual grammar of
a QR code, which is funny because the QR code is perhaps the most
aggressively contemporary object imaginable, a bureaucratic little
machine that says YOU MAY NOW LEAVE THIS PHYSICAL REALITY THROUGH YOUR
TELEPHONE, and here it is sitting among funerary stones.

Which is not subtle.

Good.

The code is not there to explain the photograph. We are trying very hard
not to explain the photograph, including right now, in a project
description whose institutional purpose is unfortunately to explain the
project.

This is the first contradiction.

Keep it.

The code opens another channel. A voice. A field recording. A sound made
in the place. A person speaking back from inside an image in which they
previously appeared to have no power except to be looked at. Then
perhaps another person. Then another practice. Photography leaks.

Image → voice → sound → person → practice → relation → next.

This chain looks suspiciously neat when written like that.

Reality will presumably damage it.

That is part of the mechanism.

The word \emph{ecclesia} matters here because before the church became
an architectural object, institutional structure, hierarchy, theology,
property portfolio and occasionally very impressive roof,
\emph{ekklesía} named an assembly. People called together. The people,
not the building.

The abandoned church therefore presents an almost absurdly literal
inversion: the architecture survives and the assembly has vanished.

TAKE OVER proposes to run the operation backward.

Not restoration.

Absolutely not restoration.

Nobody is rebuilding the lost congregation, agreeing on a doctrine,
electing a supreme leader, or politely reproducing the social order
whose remains are currently making such good photographs.

Instead:

remains → call → response → relation → assembly?

Question mark mandatory.

Because community is one of those words that starts dying the instant
everybody agrees they are doing it.

A community is not a picture of several people standing near one
another. It is not the word COMMUNITY printed in 48-point type at the
entrance of an institution. It is not consensus. It might not even be
pleasant.

It is a set of consequences produced by people becoming mutually
available to one another.

Which brings us to the arithmetic.

The call permits fifteen photographs.

We have more than fifteen photographs.

So naturally we submit twelve.

\[
12+\square+\square+\square=15
\]

Three positions remain unallocated.

This is not minimalism.

It is not curatorial restraint.

It is not because somebody forgot to export three TIFFs.

The empty positions are occupied by the possibility that somebody not
presently inside the project might alter it.

Twenty percent of scarce exhibition territory is therefore sacrificed to
a person who does not exist yet, at least not for us.

Which is economically idiotic.

Which is why it means something.

The initial participants remain identifiable. KUMIORI brings one
photographic trajectory. Ave Palm enters with another, abandoned
architectures, bodies, movement, cyanotype. Mai-Brit Tänava is
photographed, then speaks, which changes her position inside the
apparatus from represented body to speaking agent. Kenn-Eerik Kannike
introduces situated sound, electronic performance, listening.

Nobody gets dissolved into a soothing collective pronoun.

``We'' is dangerous here.

We will use it anyway.

But carefully.

The work is not the sum of these contributions. The sum is trivial. Four
people plus photographs plus audio plus a website does not automatically
equal community any more than putting twelve strangers in a lift
produces a political movement.

The work is what begins occurring between the elements.

\[
\mathcal W=\{\text{image, voice, sound, body, practice}\}-\mathcal R
\]

where $\mathcal R$ is the unstable part, the part nobody owns, the part
that grows edges.

The exhibition begins on a wall because the wall is what Fotografiska
gives us, but very quickly the wall becomes a problem.

Good.

Take over the wall.

Then go through it.

A photograph opens a sound. The sound opens a person. The person
introduces a practice. A practice may become movement. Movement might
become performance. Performance might require food. Food requires
tables. Tables produce conversation. Conversation produces another
person. Somebody brings a record. Somebody brings a projector. Somebody
says this is ridiculous.

Excellent.

The exhibition has started working.

The opening therefore cannot simply be the ceremonial moment at which
completed objects are admired beside beverages.

The opening is an activation.

Photography + voice + situated sound + movement.

After that, the sentence becomes conditional.

Maybe concert.

Maybe performance.

Maybe food.

Maybe night.

Maybe nothing we currently know how to name in the PDF, which is
slightly inconvenient for a PDF but useful for an artwork.

The abandoned church remains underneath all of this as a warning.

Assemblies become institutions.

Institutions become architecture.

Architecture survives its assembly.

Then somebody finds the ruins.

We would like to interrupt that cycle somewhere before the part where
the living thing becomes a mausoleum of itself.

So TAKE OVER does not ask how to build a permanent community.

It asks something more unstable:

How can an assembly form, act, alter its environment, and then pass
something onward before becoming the thing that must later be taken
over?

Take over the wall.

Take over the sound.

Take over the opening.

Take over the mechanism.

Then, crucially:

\textbf{PASS IT ON.}
